% 墨经原生书籍类的项目设置：A4打印，五号正文，单栏。
\geometry{a4paper,textwidth=172mm,textheight=256mm,inner=19mm,top=21mm,
  headheight=14pt,headsep=4.5mm,footskip=7.5mm}
\setlength{\parskip}{0.15em}
\xeCJKsetcharclass{"25A0}{"25FF}{1}
\xeCJKsetcharclass{"00B7}{"00B7}{0}
\xeCJKsetcharclass{"2013}{"2013}{0}
\xeCJKsetcharclass{"2103}{"2103}{1}
\setcounter{secnumdepth}{2}
\makeatletter
\renewcommand{\@pnumwidth}{2.2em}
\renewcommand{\@tocrmarg}{3.2em}
\makeatother
% A4讲义的标题留白：保留墨经的章序数、字体和细线。
\ctexset{
  chapter={beforeskip=8mm,afterskip=6mm},
  section={beforeskip=.5\baselineskip plus .2\baselineskip,
           afterskip=.25\baselineskip},
  subsection={beforeskip=.3\baselineskip plus .15\baselineskip,
              afterskip=.2\baselineskip}
}
\booktitle{高中化学系统讲义}
\booksubtitle{选择性必修1、2、3与高考综合}
\bookedition{面向广东新高考}
\publisher{}{}{2026}
